\documentclass[../main.tex]{subfiles}
\graphicspath{{\subfix{../Images/}}}

\begin{document}
\onecolumn

% =====================================================================
% Reconstructed thesis core
% Self-Organising Warehouse Systems:
% Closed-loop coupling of storage assignment and MAPD
%
% This file is a chapter fragment. Adapt citation keys to the main .bib file.
% Required packages: amsmath, amssymb, mathtools, booktabs, tabularx
% =====================================================================

\section{Introduction}
\label{sec:introduction}

Automated warehouses combine two tightly related decision layers. The storage
layer determines where stock-keeping units (SKUs) are placed, whereas the
transport layer assigns pickup--delivery tasks to mobile agents and routes the
agents through shared infrastructure. The former is commonly studied as the
Storage Location Assignment Problem (SLAP), including dynamic variants in
which assignments change over time. The latter is commonly studied through
Multi-Agent Path Finding (MAPF) and Multi-Agent Pickup and Delivery (MAPD),
where multiple agents repeatedly execute tasks while avoiding collisions
\cite{ma2017lifelong,salzman2020research,silva2020integrating}.

These layers interact through the spatial demand they induce. Placing a
frequently requested SKU at a storage location changes which vertices are used
as pickup locations and which edges carry traffic. Conversely, congestion and
waiting time change the effective cost of accessing a storage location, even
when that location is close to a delivery station in shortest-path distance.
Optimising either layer while treating the other as fixed can therefore produce
a configuration that is locally attractive but poor at system level.

This thesis investigates that interaction as a closed-loop control problem on
a common warehouse graph. A storage policy periodically updates the placement
of SKUs using observed demand and traffic statistics. Between storage updates,
a MAPD controller assigns and routes agents for the resulting stream of tasks.
The study compares a fixed-storage baseline with closed-loop variants under
centralised, section-based, and decentralised controllers. The comparison
focuses on task service time, throughput, movement cost, congestion, traffic
concentration, and computational cost.

The thesis is a simulation-based study of homogeneous mobile agents. Human
workers, physical inventory relocation, robot kinematics, battery constraints,
and uncertainty in execution are outside the primary model unless introduced
explicitly as extensions. In particular, a storage update changes the
assignment used to generate future pickup tasks; the cost and physical process
of relocating inventory are not silently assumed to be free. This abstraction
keeps the central research question feasible while making its limitations
explicit.

The main contribution is not a new general-purpose MAPF, MAPD, or SLAP solver.
It is (i) a consistent mathematical model of the coupled storage--routing
system, (ii) an implementation in which both layers use the same graph and
observed traffic state, and (iii) a controlled empirical analysis of when
feedback improves performance and when it amplifies congestion or concentrates
traffic on a small part of the warehouse.

\section{Project Description}
\label{sec:project-description}

\subsection{Aim and scope}

The project aims to determine how feedback between storage assignment and
multi-agent pickup-and-delivery control affects warehouse performance. The
warehouse is represented by a finite graph. Orders generate pickup--delivery
tasks whose pickup locations depend on the current storage configuration. A
fleet of agents serves these tasks online, and measurements from their movement
are subsequently used by the storage policy.

The independent experimental variables are:
\begin{enumerate}
  \item storage mode: fixed or periodically updated;
  \item controller architecture: centralised, section-based, or decentralised;
  \item congestion sensitivity: disabled or enabled in the storage and routing
        layers;
  \item communication: absent or local for the decentralised controller; and
  \item system load: number of agents and task-arrival rate.
\end{enumerate}
The dependent variables are mean service time, throughput, movement cost per
completed task, waiting caused by contention, traffic concentration, backlog,
and controller runtime.

\subsection{Research questions}

The main research question is:

\begin{quote}
\textbf{RQ:} How does closed-loop coupling between storage assignment and MAPD
control affect efficiency, congestion, traffic concentration, and robustness
in a graph-based warehouse?
\end{quote}

It is decomposed into four primary questions:
\begin{description}
  \item[RQ1:] Does closed-loop storage adaptation improve service time,
  throughput, or movement cost relative to fixed storage under otherwise
  matched conditions?
  \item[RQ2:] How does the effect of coupling differ among centralised,
  section-based, and decentralised controllers?
  \item[RQ3:] Do congestion-sensitive storage and routing policies reduce
  waiting and traffic concentration, and what efficiency trade-offs do they
  introduce?
  \item[RQ4:] How do these effects change as the number of agents and the task
  arrival rate increase?
\end{description}

Local communication is treated as an ablation within the decentralised
architecture rather than as a separate research topic. This keeps the scope
centred on storage--routing coupling while still testing whether communication
moderates congestion under partial observability.

\subsection{Expected outcome}

The project will produce an experimental comparison rather than claim a
universally optimal controller. Each conclusion will be conditional on the
warehouse graphs, demand processes, controller implementations, and parameter
ranges used in the experiments. The intended outcome is a map of operating
conditions under which coupling is beneficial, neutral, or destabilising.

\section{Notation}
\label{sec:notation}

The symbol \(i\) always indexes agents and \(m\) always denotes the number of
agents. Routes are denoted by \(p_i\), while \(\pi\) is reserved for a control
policy. Sets and spaces use calligraphic or uppercase letters; individual
objects use lowercase letters.

\begin{table}[htbp]
\centering
\caption{Core notation used throughout the formulation.}
\label{tab:notation}
\small
\begin{tabularx}{\textwidth}{@{}lX@{}}
\toprule
Symbol & Meaning \\
\midrule
\(G=(V,E)\) & Directed or undirected warehouse graph. \\
\(V_{\mathrm{mov}}\subseteq V\) & Vertices agents may occupy. \\
\(V_{\mathrm{str}}\subseteq V\) & Storage locations associated with pickup tasks. \\
\(V_{\mathrm{del}}\subseteq V\) & Delivery or drop-off vertices. \\
\(V_{\mathrm{ep}}\subseteq V_{\mathrm{mov}}\) & Endpoints at which an idle agent may remain without blocking transit. \\
\(c:E\rightarrow\mathbb{R}_{>0}\) & Edge traversal cost. \\
\(d_G(v,w)\) & Shortest-path cost from \(v\) to \(w\). \\
\(\mathcal{A}=\{a_1,\ldots,a_m\}\) & Agent set; \(m=|\mathcal{A}|\). \\
\(\ell_i(t)\in V_{\mathrm{mov}}\) & Location of agent \(a_i\) at time \(t\). \\
\(u_i(t)\in\mathcal{U}(\ell_i(t))\) & Action selected by agent \(a_i\). \\
\(p_i=(\ell_i(0),\ldots,\ell_i(T))\) & Realised route of agent \(a_i\). \\
\(\mathcal{K}\) & Set of SKU types. \\
\(x_t(k,v)\in\mathbb{N}_0\) & Units of SKU \(k\) assigned to storage location \(v\) at time \(t\). \\
\(\mathcal{X}\) & Set of feasible storage configurations. \\
\(\mathbb{P}_{\mathrm{ord}}\) & Distribution governing order arrivals and contents. \\
\(\lambda_{\mathrm{task}}\) & Expected number of new tasks per timestep. \\
\(\tau_j=(r_j,s_j,g_j)\) & Task \(j\), with release time, pickup vertex, and delivery vertex. \\
\(\mathcal{Q}_t\) & Released tasks waiting for assignment or pickup at time \(t\). \\
\(\mathcal{B}_t\) & Tasks being executed at time \(t\). \\
\(\mathcal{C}_T\) & Tasks completed by horizon \(T\). \\
\(\pi\in\Pi\) & Routing and task-assignment controller. \\
\(\pi_\theta\) & Parameterised decentralised policy. \\
\(\Delta\in\mathbb{N}_{>0}\) & Storage-update interval. \\
\(\mu_t(e)\) & Measured traversal count or rate on edge \(e\). \\
\(w_t(e)\) & Measured waiting attributed to edge or its downstream contention point. \\
\(C_T\in[0,1]\) & Normalised traffic-concentration index; larger means more concentrated. \\
\bottomrule
\end{tabularx}
\end{table}

\section{Problem Formulation}
\label{sec:problem-formulation}

\subsection{Warehouse graph}

Let \(G=(V,E)\) be a finite graph. Agents occupy vertices in
\(V_{\mathrm{mov}}\subseteq V\), and an action is either waiting or traversing
an incident edge. For a directed graph,
\[
\mathcal{N}^{+}(v)=\{w\in V_{\mathrm{mov}}:(v,w)\in E\},
\qquad
\mathcal{U}(v)=\{\mathsf{wait}\}\cup
\{\mathsf{move}(w):w\in\mathcal{N}^{+}(v)\}.
\]
A wait action has cost \(c_{\mathrm{wait}}\geq 0\); a move from \(v\) to \(w\)
has cost \(c(v,w)\). Thus the realised one-step cost is
\[
\widetilde c(v,w)=
\begin{cases}
 c_{\mathrm{wait}}, & v=w,\\
 c(v,w), & (v,w)\in E.
\end{cases}
\]
This definition avoids applying the edge-cost function to a wait transition
that is not an element of \(E\).

If the experiments rely on MAPD well-formedness, the following is imposed as an
instance assumption rather than presented as a consequence of connectivity:
(i) at least \(m\) non-task endpoints are available, and (ii) between any pair
of endpoints there is a path that does not pass through another endpoint
\cite{ma2017lifelong}. If the implemented controllers do not require this
assumption, it is reported only for the baselines to which it applies.

\subsection{Agents and collisions}

Let \(\mathcal{A}=\{a_1,\ldots,a_m\}\). For every \(i\in\{1,\ldots,m\}\) and
\(t\in\{0,\ldots,T-1\}\), the transition must satisfy
\[
\ell_i(t+1)=\ell_i(t)
\quad\lor\quad
(\ell_i(t),\ell_i(t+1))\in E.
\]
A joint transition is collision-free when it satisfies both
\[
\ell_i(t+1)\neq \ell_j(t+1)
\qquad \forall i\neq j,
\]
and
\[
\neg\bigl(
\ell_i(t)=\ell_j(t+1)\land
\ell_j(t)=\ell_i(t+1)
\bigr)
\qquad \forall i\neq j.
\]
The implementation must state how infeasible simultaneous actions are handled.
Either the policy is restricted to jointly feasible actions, or a deterministic
collision-resolution operator maps proposed actions to an executed joint
action. A collision penalty alone does not establish collision freedom.

\subsection{Storage configuration}

A single-valued mapping from SKU type to storage vertex cannot represent split
inventory, quantities, or vertex capacity consistently. The storage state is
therefore defined as
\[
x_t:\mathcal{K}\times V_{\mathrm{str}}\rightarrow\mathbb{N}_0,
\]
where \(x_t(k,v)\) is the number of units of SKU \(k\) assigned to location
\(v\). Let \(b_k>0\) denote the capacity consumed by one unit of SKU \(k\), and
let \(\operatorname{cap}(v)>0\) be the capacity of \(v\). The feasible set is
\[
\mathcal{X}=
\left\{
 x\in\mathbb{N}_0^{|\mathcal{K}|\times|V_{\mathrm{str}}|}:
 \sum_{k\in\mathcal{K}}b_kx(k,v)\leq\operatorname{cap}(v)
 \quad\forall v\in V_{\mathrm{str}}
\right\}.
\]
If each SKU is intentionally restricted to one location and each location is
measured in SKU slots, this may be simplified to a mapping
\(\sigma_t:\mathcal{K}\rightarrow V_{\mathrm{str}}\), but that assumption must
then be stated explicitly.

For the abstract relocation model used here, the update at epochs
\(t\in\{\Delta,2\Delta,\ldots\}\) is
\[
x_t=F(x_{t-\Delta},\widehat\rho_t,\widehat\mu_t,\widehat w_t)
\in\mathcal{X},
\]
where the estimates are computed only from data available before the update.
The fixed-storage baseline uses the identity update
\(F_{\mathrm{fix}}(x,\cdot)=x\). It is clearer to define this baseline directly
than to use the informal value \(\Delta=\infty\).

A congestion-aware assignment score may be used, but capacity couples the SKU
decisions. Hence the update is a joint constrained optimisation,
\[
x_t\in\arg\min_{x\in\mathcal{X}}
\left[
 \alpha D(x;\widehat\rho_t)
 +\beta K(x;\widehat\mu_t,\widehat w_t)
 +\chi R(x,x_{t-\Delta})
\right],
\]
where \(D\) measures demand-weighted access distance, \(K\) measures predicted
congestion exposure, and \(R\) penalises reassignment. The term \(R\) may be
set to zero in the abstract model, but doing so must be reported because it
removes relocation cost from the optimisation.

\subsection{Orders and tasks}

At each timestep, orders are sampled from
\(\mathbb{P}_{\mathrm{ord}}\). An order is a finite multiset of SKU--quantity
pairs. It is transformed into one or more tasks according to an explicitly
specified task-generation rule. A generated task is
\[
\tau_j=(r_j,s_j,g_j),
\qquad
r_j\in\mathbb{N}_0,
\quad s_j\in V_{\mathrm{str}},
\quad g_j\in V_{\mathrm{del}}.
\]
The pickup \(s_j\) must be selected from a location with available inventory,
that is, \(x_{r_j}(k_j,s_j)>0\) for the requested SKU \(k_j\). If inventory
depletion and replenishment are not simulated, \(x_t\) is interpreted as a
placement template rather than a physical inventory state.

Task lifecycle is represented by disjoint sets: waiting tasks
\(\mathcal{Q}_t\), active tasks \(\mathcal{B}_t\), and completed tasks
\(\mathcal{C}_t\). A task enters \(\mathcal{Q}_t\) at release, moves to
\(\mathcal{B}_t\) when execution begins under the chosen convention, and enters
\(\mathcal{C}_t\) on delivery. The thesis must use one convention consistently:
either execution starts at assignment or it starts when the assigned agent
reaches the pickup vertex.

If \(d_j\) is the completion time, the service time is
\[
\zeta_j=d_j-r_j.
\]
The assignment component of controller \(\pi\) maps available agents and
waiting tasks to agent--task pairs. The routing component maps the resulting
agent targets and system information to actions. Treating both components as
part of \(\pi\) avoids describing a path planner as though it also solves task
assignment when it does not.

\subsection{Controller architectures}

The experiments compare three implemented architectures, denoted
\(\pi_{\mathrm{cen}}\), \(\pi_{\mathrm{sec}}\), and \(\pi_{\mathrm{dec}}\).
These symbols identify concrete controllers, not unrestricted classes whose
members may have incomparable capabilities.

\paragraph{Centralised controller.}
The controller observes the global state, performs task assignment, and
produces coordinated actions or collision-free paths for all agents. The
replanning trigger and planning horizon are implementation choices and must be
reported; ``centralised'' does not imply replanning at every timestep.

\paragraph{Section-based controller.}
Let \(\mathcal{Y}=\{Y_1,\ldots,Y_p\}\) be a partition of
\(V_{\mathrm{mov}}\):
\[
\bigcup_{q=1}^{p}Y_q=V_{\mathrm{mov}},
\qquad Y_q\cap Y_{q'}=\varnothing\quad(q\neq q').
\]
The controller decomposes assignment or routing by section and uses an
explicit protocol for cross-section tasks and boundary contention. Without
such a protocol, ``section-based'' is an architectural label rather than a
complete algorithm.

\paragraph{Decentralised controller.}
Each agent selects an action from a local observation using a shared policy,
\[
u_i(t)\sim\pi_\theta(\cdot\mid o_i(t)).
\]
Parameter sharing is appropriate for homogeneous agents but does not make
execution centralised. If centralised information is used during training but
not execution, the method is described as centralised training with
decentralised execution.

\subsection{Local observations and communication}

For observation depth \(d\), define
\[
V_i^{(d)}(t)=\{v\in V_{\mathrm{mov}}:d_G(\ell_i(t),v)\leq d\},
\qquad
G_i^{(d)}(t)=G[V_i^{(d)}(t)].
\]
If the current target of agent \(a_i\) is \(q_i(t)\), each visible vertex may
be annotated by
\[
\eta_i(v,t)=d_G(v,q_i(t)).
\]
All-pairs or target-specific shortest-path distances may be precomputed, but
\(\eta_i(v,t)\) itself changes when the target changes.

The communication graph is
\[
G^{\mathcal A}_t=(\mathcal{A},E^{\mathcal A}_t),
\qquad
(a_i,a_j)\in E^{\mathcal A}_t
\iff i\neq j\land
d_G(\ell_i(t),\ell_j(t))\leq r_{\mathrm{com}}.
\]
The environment graph and communication graph have different node types and
must not be conflated. A two-stage encoder can first encode the local
warehouse subgraph,
\[
z_i(t)=\operatorname{Enc}_{G}(G_i^{(d)}(t),\text{node/edge features}),
\]
and then exchange agent-level messages,
\[
h_i^{(q+1)}(t)=\phi_q\!\left(
 h_i^{(q)}(t),
 \operatorname{AGG}\{\psi_q(h_i^{(q)}(t),h_j^{(q)}(t)):
 (a_i,a_j)\in E^{\mathcal A}_t\}
\right),
\]
with \(h_i^{(0)}(t)=z_i(t)\). This matches the verbal claim that both graph
structures are processed; a message-passing equation over agents alone does
not encode the local warehouse graph.

\subsection{Stochastic-game model}

The decentralised learning problem is represented as a partially observable
stochastic game
\[
\mathcal{M}=\left(
 \mathcal{A},\mathcal{S},
 \{\mathcal{U}_i\}_{i=1}^{m},
 P,\{\mathcal{O}_i\}_{i=1}^{m},
 \{O_i\}_{i=1}^{m},
 \{R_i\}_{i=1}^{m},\gamma
\right),
\]
where \(\mathcal{S}\) is the state space, \(P\) is the transition kernel,
\(O_i\) is the observation function, and
\(R_i:\mathcal{S}\times\mathcal{U}^m\times\mathcal{S}\to\mathbb{R}\) is the
reward function. This distinguishes spaces and functions from realised states,
observations, and rewards.

A local congestion feature can be defined without counting the focal agent as
\[
\delta_i(t)=
\frac{
 |\{a_j\in\mathcal{A}\setminus\{a_i\}:
 \ell_j(t)\in V_i^{(r_{\mathrm{cng}})}(t)\}|
}{
 \max\{1,|V_i^{(r_{\mathrm{cng}})}(t)|-1\}
}.
\]
The reward may combine progress, task completion, waiting, invalid proposals,
and congestion. For a shared cooperative policy, the learning objective is
stated at team level,
\[
\max_\theta J(\theta)
=
\max_\theta
\mathbb{E}_{\pi_\theta}
\left[
\sum_{t=0}^{T-1}\gamma^t
\frac{1}{m}\sum_{i=1}^{m}R_i(t)
\right].
\]
This is consistent with one shared parameter vector; maximising an unspecified
single agent's return is not.

\subsection{Coupled dynamics}

Let \(s_t\in\mathcal{S}\) include agent states, task lifecycle sets, and the
current storage configuration. For non-update timesteps, the fast dynamics are
\[
s_{t+1}\sim P(\cdot\mid s_t,\mathbf{u}_t),
\qquad
\mathbf{u}_t\sim\pi(\cdot\mid s_t)
\]
for a centralised controller, with the corresponding product of local policies
for decentralised execution. At storage-update epochs, the slow dynamics also
apply \(F\) to determine the configuration used by future task generation.
Thus the storage policy changes the routing problem through the distribution of
future pickup vertices, while realised routing statistics change later storage
updates.

\subsection{Evaluation measures}

For completed tasks \(\mathcal{C}_T\), define throughput and mean service time
as
\[
\Lambda_T=\frac{|\mathcal{C}_T|}{T},
\qquad
\overline\zeta_T=
\frac{1}{|\mathcal{C}_T|}
\sum_{\tau_j\in\mathcal{C}_T}(d_j-r_j).
\]
The movement cost per completed task is
\[
\overline c_T=
\frac{
\sum_{i=1}^{m}\sum_{t=0}^{T-1}
\widetilde c(\ell_i(t),\ell_i(t+1))
}{|\mathcal{C}_T|}.
\]
Runs with \(|\mathcal{C}_T|=0\) are reported separately rather than assigned an
arbitrary ratio. Because completion-only service time can hide an increasing
backlog, \(|\mathcal{Q}_T|+|\mathcal{B}_T|\) and the age of unfinished tasks are
reported alongside it.

Let \(\mu_T(e)\) be the number of traversals of edge \(e\) during the evaluation
window and let
\(p_T(e)=\mu_T(e)/\sum_{e'\in E}\mu_T(e')\). Normalised entropy is
\[
H_T=-\frac{1}{\log |E^+_T|}
\sum_{e\in E^+_T}p_T(e)\log p_T(e),
\qquad
E^+_T=\{e\in E:\mu_T(e)>0\},
\]
when \(|E^+_T|>1\). Since larger entropy means more dispersed traffic, define
concentration as
\[
C_T=1-H_T.
\]
This direction is essential: adding \(+H_T\) to a minimised cost would reward
concentration rather than penalise it.

The system-level evaluation functional is
\[
\mathcal{J}(F,\pi)=
q_{\mathrm{svc}}\,\mathbb{E}[\overline\zeta_T]
+q_{\mathrm{mov}}\,\mathbb{E}[\overline c_T]
+q_{\mathrm{wait}}\,\mathbb{E}[W_T]
+q_{\mathrm{conc}}\,\mathbb{E}[C_T]
+q_{\mathrm{back}}\,\mathbb{E}[B_T],
\]
where all \(q_\bullet\geq0\), \(W_T\) is a precisely implemented waiting
measure, and \(B_T=|\mathcal{Q}_T|+|\mathcal{B}_T|\). Throughput is reported as
a separate maximisation-oriented measure rather than inserted into the cost
with an ambiguous sign.

The formal decision problem is therefore
\[
(F^\star,\pi^\star)
\in
\arg\min_{F\in\mathcal{F},\,\pi\in\Pi}
\mathcal{J}(F,\pi)
\quad\text{subject to storage feasibility, valid task transitions, and
collision-free executed actions.}
\]
The empirical thesis does not need to solve this joint optimisation globally.
It estimates the effects of selected \((F,\pi)\) combinations under a factorial
experimental design.

\section{Background}
\label{sec:background}

\subsection{Storage assignment}

SLAP determines where items should be stored under constraints such as storage
capacity, compatibility, and demand. Common objectives include reducing travel
and improving space utilisation. Dynamic SLAP extends the decision over time as
demand or inventory conditions change. For this thesis, the important point is
not a catalogue of all SLAP variants but the mechanism by which placement
changes the distribution of pickup locations. The background should therefore
introduce static and dynamic assignment, demand-frequency heuristics, capacity
constraints, and relocation cost before presenting the chosen update rule
\cite{puiseau2022dynamic,silva2020integrating}.

\subsection{MAPF and MAPD}

Classical MAPF assigns each agent a start and goal on a graph and seeks
collision-free paths, commonly under makespan or sum-of-cost objectives. MAPD
extends this setting to a continuing stream of pickup--delivery tasks and adds
task assignment and repeated planning. This distinction matters because the
thesis evaluates long-run service and throughput rather than only the cost of a
single set of paths \cite{ma2017lifelong,salzman2020research}.

The background should define vertex and swap conflicts, endpoints,
well-formedness, service time, task assignment, and replanning. Solver families
should be discussed only to the depth needed to justify the selected baselines.
A long list of algorithms that are neither implemented nor used analytically
would obscure the methodological choices.

\subsection{MARL and partial observability}

A decentralised agent acts from a local observation while the environment
evolves under the joint actions of all agents. This motivates a partially
observable stochastic-game model. Relevant concepts are parameter sharing for
homogeneous agents, centralised training with decentralised execution,
non-stationarity, action masking, reward shaping, and evaluation over multiple
random seeds. Algorithmic background should focus on the learning algorithm
actually used in the experiments, with alternatives introduced only where they
explain a design choice or baseline \cite{lau2022mapd}.

\subsection{Graph representations}

The warehouse topology may be non-lattice, so a graph encoder is more natural
than a convolution tied to a rectangular grid. A fixed-depth induced subgraph
provides each agent with a bounded local observation, while shortest-path
potentials communicate target direction even when the target lies outside that
subgraph. A separate time-varying interaction graph represents which agents
can exchange information. These two graph levels must be introduced separately
because they encode different entities and relations
\cite{vanknippenberg2021time,pham2024crowd}.

\subsection{Congestion and emergence}

Congestion is an operational phenomenon and must be tied to observable
quantities such as blocked moves, waiting time, queue length, vertex occupancy,
or edge usage. Traffic concentration is related but not identical: a flow can
be highly concentrated without causing delay at low load, and dispersed flow
can still be inefficient. The thesis therefore reports congestion and
concentration as separate outcomes and uses ``commons-like'' as an
interpretive analogy only when individual or local incentives systematically
increase shared-resource overuse.

\section{Related Work}
\label{sec:related-work}

\subsection{Integrated storage and picking}

Work integrating storage location and order picking establishes that placement
and picking decisions should not always be optimised independently
\cite{silva2020integrating}. The present thesis differs by making realised
multi-agent traffic and congestion part of the feedback to later storage
updates. The relevant gap is therefore not simply that storage and routing have
never been combined, but that closed-loop, routing-level congestion feedback
has received less attention in the graph-based MAPD setting.

\subsection{Online pickup and delivery}

MAPD research formalises lifelong operation with online pickup--delivery tasks
and provides centralised or token-based methods under well-formedness
conditions \cite{ma2017lifelong}. Related lifelong MAPF work also studies
replanning and scalability in warehouse-like environments. These methods
supply definitions and potential baselines, but they generally treat task
pickup locations as exogenous rather than as outputs of an adaptive storage
layer.

\subsection{Learned decentralised routing}

Learned MAPF and MAPD controllers trade explicit global planning for fast local
inference. Work on non-lattice graphs shows how local graph observations and
distance-to-goal attributes can support a shared learned policy
\cite{vanknippenberg2021time}. Crowd-aware GNN approaches further motivate
local communication and congestion-sensitive rewards
\cite{pham2024crowd}. The thesis uses these ideas to test an interaction they
do not by themselves resolve: how a learned routing policy behaves when task
origins are changed by an adaptive storage policy.

\subsection{Positioning of this thesis}

The thesis lies at the intersection of dynamic storage assignment, lifelong
multi-agent routing, and graph-based decentralised control. Its novelty claim
should be limited to the tested combination and experimental evidence:
closed-loop storage updates informed by realised multi-agent traffic, compared
across matched controller architectures and load levels. Claims that the two
fields are ``always'' studied independently, or that the proposed question
cannot be addressed with classical methods, are too strong unless supported by
a systematic literature review.

\section{Experimental Comparison}
\label{sec:experimental-comparison}

For each warehouse graph and demand regime, evaluate matched configurations
\[
\mathcal{E}=
\mathcal{F}_{\mathrm{storage}}
\times
\mathcal{P}_{\mathrm{controller}}
\times
\mathcal{B}_{\mathrm{congestion}}
\times
\mathcal{B}_{\mathrm{communication}}
\times
\mathcal{M}_{\mathrm{agents}}
\times
\mathcal{L}_{\mathrm{task}},
\]
where communication ablations apply only to the decentralised controller. The
fixed and adaptive storage conditions must share the same graph, exogenous
order realisations, initial agent positions, controller hyperparameters, and
evaluation horizon wherever paired comparison is intended.

Report distributions or confidence intervals across independent seeds, not
only point estimates. Separate training seeds from evaluation scenarios, and
evaluate learned controllers on order streams not used for training. Claims of
robustness should be tied to predefined perturbations, such as increased agent
count, increased arrival rate, demand shift, or graph changes, rather than used
as a synonym for high average performance.
\end{document}